
%%%%%%%%%%  scan idx 335  |  folio 328  %%%%%%%%%%


S G A \quad 7


\underline{Exposé XIX}


\underline{\underline{Le théorème de Noether}}


\underline{par \quad P. Deligne}


\underline{Sommaire}

1. Enoncé des théorèmes \hfill 1

2. Complément à XVIII \hfill 4

3. Preuve de 1.3 \hfill 5

4. Deuxième complément à XVIII \hfill 8

5. Les cas d'exception au théorème de Noether. \hfill 11


\underline{1. Enoncé des théorèmes}

Le présent exposé est une mise au goût du jour de la démonstration de Lefschetz [3] page 108 du théorème de Noether (1.2) ci-dessous.

1.1 \qquad Soit $k$ un corps algébriquement clos de caractéristique $p$ . Soient $d \geq 1$ et $\underline{a} = (a_i)_{1 \leq i \leq d}$ une suite d'entiers $\geq 1$ . Une intersection complète $V$ de dimension $n$ et de multidegré $\underline{a}$ dans l'espace projectif $\mathbb{P}^{n+d}$ sur $k$ sera dite \underline{générique} si, sur une extension de $k$ , elle est projectivement isomorphe à une intersection complète définie par des équations dont les coefficients sont algébriquement indépendants sur le corps premier. Dans ce qui suit et sauf mention expresse du contraire, \underline{on suppose que} $d \geq 1$ et que $a_i \geq 2$ .

\underline{Théorème} 1.2 \qquad \underline{Soit} $V \subset \mathbb{P}^{2+d}$ \underline{une surface intersection complète générique de multidegré} $\underline{a}$ . \underline{On suppose que} $\underline{a} \neq (2), (3), (2, 2)$ . \underline{Alors, toute courbe (réduite) tracée sur} $V$ \underline{est l'intersection (schématique) de} $V$ \underline{et d'une hypersurface de} $\mathbb{P}^{2+d}$ .

Montrons que le théorème résulte de son corollaire 


%%%%%%%%%%  scan idx 336  |  folio 329  %%%%%%%%%%

\begin{equation}\tag{1.2.1}
 \mathrm{Pic}\ (\mathbb{P}^{2+d}) = \mathbb{Z} \overset{\sim}{\longrightarrow} \mathrm{Pic}(V)\ .
\end{equation}
Si $C$ est une courbe tracée sur $V$ , (1.2.1) implique que $\mathcal{O}(C) \simeq \mathcal{O}(n)$ pour $n$ convenable; si $H$ est une section hyperplane, $n = (\Pi\ a_i)^{-1}.(C, H) > 0$ . Il reste à prouver que le système linéaire induit sur $V$ par $|n\ H|$ est complet. Ecrivons $V$ comme intersection d'hypersurfaces $H_i$ de degré $a_i$ et soit $V_i = H_1 \cap \ldots \cap H_i$ . On vérifie par récurrence sur $i$ que le système linéaire induit par $|n\ H|$ sur $V_i$ est complet: dans la suite exacte

\begin{equation*}
H^0(V_i,\ \mathcal{O}(n)) \longrightarrow H^0(V_{i+1},\ \mathcal{O}(n)) \longrightarrow H^1(V_i,\ \mathcal{O}(n - a_{i+1}))\ ,
\end{equation*}
le dernier groupe est nul (FAC) .

\underline{Théorème} 1.3 \qquad \underline{Soit} $V \subset \mathbb{P}^{2n+d}$ \underline{une intersection complète de dimension paire} $2n \neq 0$ \underline{et de multidegré} $\underline{a}$ \underline{générique.} \underline{On suppose vérifiée l'une des conditions suivantes}

(a) \quad $k$ \underline{n'est pas de caractéristique} $2$ , \underline{et} $(2n; \underline{a})$ \underline{est distinct de} $(2n; 2)$, $(2n; 2, 2)$ , $(2; 3)$ .

(b) \quad \underline{le nombre de Hodge} $h^{2n,0}$ \underline{est non nul,} \underline{i.e.} $\Sigma\ (a_i-1) > 2n$ .

\underline{Alors, toute classe de cohomologie algébrique dans} $H^{2n}(V,\ \mathbb{Q}_\ell(n))$ \underline{est un multiple rationnel de la classe} $\eta^n$ \underline{des sections linéaires.}

Déduisons 1.2.1 de 1.3. Les cas d'exceptions dans 1.2 sont exactement ceux pour lesquels $\Sigma(d_i -1) \leq 2$ ; on peut donc appliquer 1.3 (b) à la surface $V$ de 1.2. On sait que $\mathrm{Pic}^{\tau}(V) = 0$ (XI 1.8); donc que $\mathrm{Pic}(V)$ s'injecte dans $H^2(V,\ \mathbb{Q}_\ell(1))$ . D'après 1.3, tout élément de $\mathrm{Pic}(V)$ est un multiple rationnel de la classe $\eta$ de $\mathcal{O}(1)$ , et on conclut par (XI 1.8).

Il est clair que la validité de 1.2 et 1.3 ne dépend pas de l'intersection complète générique choisie, mais seulement de $p, n$ et $\underline{a}$ .


%%%%%%%%%%  scan idx 337  |  folio 330  %%%%%%%%%%

\underline{Théorème 1.4.} \quad \underline{Soit} $X \subset \mathbb{P}$ \underline{une variété projective lisse purement de dimension} $2n + 1$ \underline{sur} $k$ . \underline{Soit} $\ell \subset \check{\mathbb{P}}$ \underline{un pinceau de Lefschetz de sections hyperplanes de} $X$ . \underline{Pour} $k$ \underline{de caractéristique} 2, \underline{on suppose que} $\ell$ \underline{est même un pinceau de Lefschetz générique.}

\underline{On suppose que} $X$ \underline{vérifie} (LV)  (XVIII 5.2.2).

\underline{Soit} $V$ \underline{la section hyperplane générique dans le pinceau} $\ell$ . \underline{L'une des deux conditions suivantes est vérifiée.}

(a) \underline{Aucune classe de cohomologie} $\neq 0$ \underline{dans la partie évanescente} $\mathrm{Ev}(V, \mathbb{Q}_\ell(n))$ \underline{du groupe de cohomologie} $H^{2n}(V, \mathbb{Q}_\ell(n))$ \underline{n'est algébrique.}

(b) \underline{Les classes de cohomologie algébrique engendrent} $\mathrm{Ev}(V, \mathbb{Q}_\ell(n))$ .

Précisons le sens du mot générique dans 1.4.

(a) Soit $G$ la grassmanienne des droites de l'espace projectif $\check{\mathbb{P}}$ dual de $\mathbb{P}$ . Appelons ``point géométrique'' de $G$ un morphisme du spectre d'un corps algébriquement clos dans $G$ .

Tout point géométrique $\bar{s} \to G$ définit un pinceau de sections hyperplanes $\ell$ de la variété déduite de $X$ par extension des scalaires de $k$ à $k(\bar{s})$ . On dit que $\ell$ est générique si $\bar{s}$ est localisé au point générique de $G$ .

(b) Soit $\ell$ un pinceau de sections hyperplanes (une droite de $\check{\mathbb{P}}$ ). Tout point géométrique $\bar{s} \to \ell$ définit une section hyperplane de la variété déduite de $X$ par extension des scalaires de $k$ à $k(\bar{s})$ . On dit qu'elle est générique si $\bar{s}$ est localisé au point générique de $\ell$ .

Soient $\bar{s}$ comme ci-dessus, d'image dans $\ell$ le point générique $s$ et $S \subset \ell$ l'ensemble fini des points de $\ell$ donnant lieu à une section hyperplane singulière. Le groupe de monodromie $\pi_1(\ell - S, \bar{s})$ , quotient de $\mathrm{Gal}(k(\bar{s})/k(s))$ , agit sur $\mathrm{Ev}(V, \mathbb{Q}_\ell(n))$ . Si $k$ n'est pas de caractéristique 2, cette action est irréductible (XVIII 6.7); au n$^{\circ}$ 2 ci-dessous, on montrera que ce résultat reste valable sans hypothèse de caractéristique pour un pinceau générique. Le sous-espace de $\mathrm{Ev}(V, \mathbb{Q}_\ell(n))$ engendré par les classes de cohomologie algébrique est stable sous le groupe de monodromie (i.e. sous Galois); c'est donc tout ou rien, ce qui prouve 1.4.


%%%%%%%%%%  scan idx 338  |  folio 331  %%%%%%%%%%

\underline{Corollaire 1.5} \quad \underline{Soit} $V \subset \mathbb{P}^{2n+d}$ \underline{une intersection complète de dimension paire} $2n$ \underline{et de multidegré} $\underline{a}$ \underline{générique. On a soit}

(a) \underline{Toute classe de cohomologie algébrique dans} $H^{2n}(V, \mathbb{Q}_\ell(n))$ \underline{est multiple rationnel de} $\eta^n$ .

(b) $H^{2n}(V, \mathbb{Q}_\ell(n))$ \underline{est engendré par des classes de cohomologie algébrique.}

\underline{Si} (b) \underline{est vérifié, pour} toute \underline{intersection complète de dimension} $2n$ \underline{et multidegré} $\underline{a}$ \underline{sur un corps algébriquement clos de caractéristique} $p$ , $H^{2n}(V, \mathbb{Q}_\ell(n))$ \underline{est engendré par les classes algébriques.}

On applique 1.4 pour $X$ une intersection complète de dimension $2n+1$ et de multidegré $(a_1 \ldots a_{d-1})$ générique et pour $\ell$ un pinceau générique de sections de $X$ par des hypersurfaces de degré $a_d$ (cf. XVII 2.4). En fait, on pourrait aussi prendre $X$ de multidegré $(a_1 \ldots a_d)$ et $\ell$ un pinceau générique de sections hyperplanes. La remarque finale résulte de ce qu'une classe de cohomologie algébrique reste algébrique par spécialisation.


\underline{2. Complément à XVIII}

\underline{Proposition 2.1} \quad \underline{Soient} $X \subset \mathbb{P}$ \underline{une variété projective lisse purement de dimension paire} $2n$ \underline{sur un corps algébriquement clos} $k$ , \underline{et} $\ell \subset \check{\mathbb{P}}$ \underline{un pinceau de Lefschetz générique de sections hyperplanes de} $X$ . \underline{Soient} $s$ \underline{un point géométrique de} $\ell$ , $H_s$ \underline{la section hyperplane correspondante,} $S \subset \ell$ \underline{l'ensemble des valeurs exceptionnelles du pinceau et} $\pi_1(\ell - S, s)$ \underline{le groupe de monodromie.}

\underline{On suppose que} $X$ \underline{vérifie} (LV) .

\underline{Alors, la représentation de} $\pi_1(\ell - S, s)$ \underline{sur la cohomologie évanescente} $\mathrm{Ev}(H_s, \mathbb{Q}_\ell) \subset H^{2n}(H_s, \mathbb{Q}_\ell)$ \underline{est absolument irréductible.}

Rappelons que $\mathrm{Ev}(H_s, \mathbb{Q}_\ell)$ est l'orthogonal de l'image de $H^{2n}(X, \mathbb{Q}_\ell)$ dans $H^{2n}(H_{s,} \mathbb{Q}_\ell)$ , image égale au sous-groupe des invariants sous le groupe de monodromie (XVIII 5.6.10).


%%%%%%%%%%  scan idx 339  |  folio 332  %%%%%%%%%%

La droite projective $\mathbb{P}^1$ est simplement connexe. Le groupe $\pi_1(\ell - S, s)$ est donc (topologiquement) engendré par les sous-groupes d'inertie en les points de $S$ - et leurs conjugués. Son image dans $\mathrm{GL}(H^n(H_s, \mathbb{Q}_\ell))$ est donc (topologiquement) engendrée par les $x \longrightarrow x - (-1)^n(x\ \delta)\delta$ , pour $\delta$ conjugué à un cycle évanescent. L'image de $H^{2n}(X, \mathbb{Q}_\ell)$ dans $H^{2n}(H_s, \mathbb{Q}_\ell)$ est donc l'orthogonal du sous-espace engendré par les cycles évanescents et leur conjugués: $\mathrm{Ev}(H_s, \mathbb{Q}_\ell)$ est engendré par les cycles évanescents et leur conjugués.

Pour un pinceau de Lefschetz \underline{générique} , les arguments du n° 6 de XVIII montrent encore que les cycles évanescents $\delta$ relatifs aux différents points de $S$ sont conjugués entre eux.

Soit $V \subset \mathrm{Ev}(H_s, \mathbb{Q}_\ell)$ un sous-espace invariant $\neq 0$ . Il existe un conjugué $\delta$ d'un cycle évanescent et $v \in V$ tels que $(v, \delta) \neq 0$ (c'est ici qu'on utilise (LV)). Puisque $v - (-1)^m(v\ \delta)\delta \in V$ , on a $\delta \in V$ . Le sous-espace $V$ , contenant un cycle évanescent (et tous ses conjugués) coïncide avec $\mathrm{Ev}(H_s, \mathbb{Q}_\ell)$ .

Cet argument reste valable après extension du corps de coefficients $\mathbb{Q}_\ell$ , d'où l'absolue irréductibilité.

\underline{Remarque} 2.2 \quad Lorsqu'on omet l'hypothèse (LV) , les arguments qui précèdent prouvent encore que tout sous-espace invariant de $H^{2n}(H_s, \mathbb{Q}_\ell)$ contient $\mathrm{Ev}(H_s, \mathbb{Q}_\ell)$ ou est contenu dans l'image de $H^{2n}(X, \mathbb{Q}_\ell)$ . Le même résultat vaut pour $X$ de dimension paire (et tout pinceau de Lefschetz).


\underline{3. Preuve de 1.3}

3.1 \quad Il s'agit de prouver que, sous les hypothèses du théorème, 1.5 (b) n'est pas vérifié. Le cas où $h^{0,2n} \neq 0$ sera traité par Katz en XXI 4.1 (la condition $h^{0,2n} \neq 0$ équivaut à $\Sigma\, d_i - (2n + r + 1) \geq 0$ , ce nombre étant l'entier $m$ tel que $\Omega^{2n}_V \simeq \mathcal{O}(m)$ ). Les résultats du présent exposé ne servirons plus dans la suite de séminaire, de sorte qu'il n'y a pas de cercle vicieux.


%%%%%%%%%%  scan idx 340  |  folio 333  %%%%%%%%%%

3.2 \quad Lorsque $k = \mathbb{C}$ , on démontre 1.3 en notant que si 1.5 (b) est vérifié, la cohomologie de $V$ est purement de type $(n,n)$ , la cohomologie algébrique étant de type $(n,n)$ . Ceci n'est le cas que dans les cas exceptionnels exclus de 1.3 (XI 2.9). Par le principe de Lefschetz, 1.3 est donc vrai en caractérisitque $0$ .

\underline{Proposition} 3.3 \quad \underline{Plaçons-nous sous les hypothèses de} 1.4, \underline{et supposons} (b) \underline{vérifié.} \underline{Soit} $\mathrm{Ev}(V, \mathbb{Q})$ \underline{le} $\mathbb{Q}$\underline{-sous-espace vectoriel de} $\mathrm{Ev}(V, \mathbb{Q}_\ell(n))$ \underline{engendré par les classes de cohomologie algébrique.}

(a) \underline{On a} $\mathrm{Ev}(V, \mathbb{Q}) \otimes \mathbb{Q}_\ell \stackrel{\sim}{\longrightarrow} \mathrm{Ev}(V, \mathbb{Q}_\ell(n))$ . \underline{Sur} $\mathrm{Ev}(V, \mathbb{Q})$ , \underline{la forme d'intersection est rationnelle définie, de signe} $(-1)^n$ .

(b) \underline{Les cycles évanescents et leur conjugués sont dans} $\mathrm{Ev}(V, \mathbb{Q})$ ; \underline{ils y forment un système de racine irréductible de type} A, D \underline{ou} E .

(c) \underline{Le groupe de monodromie (image de} $\pi_1(\ell - S, \bar{s})$ \underline{dans} $\mathrm{GL}(\mathrm{Ev}(V, \mathbb{Q}))$ ) \underline{est fini; c'est le groupe de Weil du système de racines} (b) .

Sauf le caractère défini de la forme d'intersection, (a) résulte de ce que la forme d'intersection est non dégénérée sur $\mathrm{Ev}(V, \mathbb{Q}_\ell(n))$ (par (LV)) et rationnelle sur $\mathrm{Ev}(V, \mathbb{Q})$ .

La condition 1.4 (b) équivaut à dire qu'il existe sur $V$ des cycles algébriques modulo équivalence numérique non invariants par la monodromie. Cette condition est donc indépendante de $\ell$ . Soit $x$ un tel cycle algébrique, et $\delta$ le cycle évanescent relatif à $s_o \in S$ et à un chemin ch de $\bar{s}$ (définissant $V$ ) à $s_o$ . Pour ch convenable, $(x\ \delta) \neq 0$ . Le transformé galoisien $\sigma(x) = x-(-1)^n(x\delta)\delta$ de $x$ est encore algébrique; $(x\delta)\delta$ est donc algébrique, image d'un cycle $y$ indépendant de $\ell$ . Le nombre rationnel $\frac{(-1)^n}{2}(y, y) = (x\delta)^2$ est donc indépendant de $\ell$ ; c'est un carré dans tous les corps $\mathbb{Q}_\ell$ , donc un carré, et $\pm (x, \delta)$ est dans $\mathbb{Q}$ et indépendant de $\ell$ . Ceci montre non seulement que $\delta \in \mathrm{Ev}(V, \mathbb{Q})$ , mais encore que le cycle algébrique modulo équivalence numérique $\pm\, \delta = (\sigma(x) -x)/(x, \delta)$ , est indépendant de $\ell$ .


%%%%%%%%%%  scan idx 341  |  folio 334  %%%%%%%%%%

Un cycle algébrique sur une variété $V$ est toujours défini sur une extension finie de tout corps de définition de $V$ . Il en résulte que $\pi_1(\ell - S, s)$ agit sur $\mathrm{Ev}(V, \mathbb{Q})$ via un groupe fini. Ce groupe est engendré par les réflexions associées aux cycles évanescents $\delta$ et leur conjugués (XVIII 6.6.1 et n$^\circ$ 2). Les $\delta$ et leur conjugués forment donc un système de racines de groupe de Weil le groupe de monodromie. De plus, toutes les racines de ce système de racines sont conjuguées sous le groupe de Weil: il est de type $A$, $D$ ou $E$ .

La forme d'intersection est, à un facteur près, l'unique forme invariante sous le groupe de Weil. Elle est donc définie. Son signe se détermine en calculant

\begin{equation*}
(\delta, \delta) = (-1)^n . 2 \quad .
\end{equation*}
\underline{Proposition} 3.4 \quad \underline{Si} $k = \mathbb{C}$ , \underline{et sous les hypothèses de} 1.4, \underline{les conditions suivantes sont équivalentes}

(i) \underline{le groupe de monodromie est fini};

(ii) \underline{la cohomologie évanescente} $\mathrm{Ev}(V)$ \underline{est purement de type} $(n, n)$ .

Soit $\widetilde{X} \longrightarrow \ell$ le fibré de fibres les sections hyperplanes appartenant au pinceau $\ell$ , et $V_s$ la fibre en $s \in \ell(\mathbb{C})$ . Les $\mathrm{Ev}(V_s)$ forment une variation algébrique de structures de Hodge sur $\ell$ . D'après le théorème de la partie fixe [2] 4.1.2 appliqué à un revêtement fini de $\ell$ (cf. [2] 4.1.3.3), (i) implique que la structure de Hodge des $\mathrm{Ev}(V_s)$ est localement constante. Le groupe de monodromie agit donc par automorphismes de la structure de Hodge. Pour $T$ une transformation de Picard-Lefschetz, $\mathrm{Im}(T-1)$ est une sous-structure de Hodge de rang un, donc de type $(n, n)$ , engendrée par un cycle évanescent $\delta$ . Les cycles évanescents étant de type $(n, n)$ , $\mathrm{Ev}(V_s)$ tout entier l'est.

Réciproquement, si $\mathrm{Ev}(V)$ est purement de type $(n, n)$ , la forme d'intersection sur $\mathrm{Ev}(V)$ est définie et le groupe de monodromie fini (cf. [2] 4.1.3.3).


%%%%%%%%%%  scan idx 342  |  folio 335  %%%%%%%%%%

\underline{Remarque} 3.5 \quad Si les conditions équivalentes de 3.3 sont vérifiées, on voit comme en 3.3 que les cycles évanescents (et leurs conjugués) forment un système de racines de type $A$, $D$ ou $E$ dont le groupe de monodromie est le groupe de Weil.

3.6. \quad Prouvons 1.3 pour $k$ de caractéristique finie $\neq 2$ . Soient $S = \mathrm{Spec}(V)$ le spectre d'un anneau de valuation discrète d'inégale caractéristique de corps résiduel $k$ , et $X \subset \mathbb{P}^{2n+d}_S$ une intersection complète de dimension relative $2n+1$ et de multidegré $(a_1 \ldots a_{d-1})$ sur $S$ . Soit $\ell$ un pinceau de Lefschetz de sections hypersurfaces de degré $a_d$ de $X$ (en fibre spéciale, $\ell_s$ est un pinceau de Lefschetz pour $X_s$ ). Soit $f : \widetilde{X} \longrightarrow \ell \simeq \mathbb{P}^1_S$ le fibré de fibres les sections du pinceau. Le lieu exceptionnel $T \subset \ell$ est étale sur $S$ (XVII 6; on utilise ici que $p \neq 2$ ). Le faisceau $R^{2n} f_* \mathbb{Q}_\ell$ est localement constant sur $\ell - T$ ; il est donc modérément ramifié et sa monodromie est la même en caractéristique $p$ et en caractéristique $0$. Si 1.5 (b) est vérifié, on déduit donc de 3.3 (c) et 3.4 que le groupe de cohomologie $H^{2n}(V)$ des intersections complètes \underline{complexes} de dimension $2n$ et multidegré \underline{$a$} est purement de type $(n, n)$ , ce qui n'arrive que dans les cas exclus par 1.3 (XI 2.9).


\underline{4. Deuxième complément à XVIII}

Dans ce numéro, nous travaillons sur $\mathbb{C}$ et en cohomologie entière (transcendante). Nous reprenons, avec des méthodes plus proches de celles de Lefschetz, certains des résultats de XVIII. Dans XVIII, ils n'étaient prouvés qu'après tensorisation avec $\mathbb{Q}$ (on plutôt $\mathbb{Q}_\ell$ , ce qui revient au même).

4.1 \quad Soient $X \subset \mathbb{P}$ une variété projective complexe lisse purement de dimension $n + 1$ et $\ell$ un pinceau de Lefschetz de sections hyperplanes de $X$ , d'axe $A$ . On désigne aussi par $\ell$ l'espace des paramètres du pinceau, isomorphe à $\mathbb{P}^1$ . Soit $S \subset \ell$ l'ensemble des valeurs exceptionnelles. Pour $s \in \ell$ , on notera $H_s$ la section hyperplane correspondante de $X$ (de dimension $n$).


%%%%%%%%%%  scan idx 343  |  folio 336  %%%%%%%%%%

Soit $T(A)$ un voisinage tubulaire fermé de $A \cap X$ dans $X$ , de bord transverse à tous les $H_s$ .

Soient $0$, $\infty$ deux points dans $\ell - S$ , et soit $\alpha$ un difféomorphisme orienté de $\ell - S$ vers $\mathbb{P}^1(\mathbb{C})$ envoyant $0$ sur $0$ , $\infty$ sur $\infty$ et $S$ dans la circonférence de rayon un centrée en $0$ . Le composé $\alpha f$ de $\alpha$ et de la projection sur $\ell$ fait de $X - T(A)^o$ un fibré singulier en variétés à bord sur $\mathbb{P}^1(\mathbb{C})$ (une singularité au-dessus de chaque point de $\alpha(S)$ ). Nous noterons $c_s$ le segment de droite joignant $0$ à $\alpha(s)$ .

Omettant la mention de $\alpha$ , on a la figure

\begin{center}
\begin{tikzcd}[row sep=huge, column sep=huge]
{} & s_2 & s_1 \\
\cdots & 0 \arrow[u, dash, "{c_2}"'] \arrow[ur, dash, "{c_1}"'] \arrow[ul, dash]
\end{tikzcd}
\end{center}
% STRUCTURE: Freehand line figure (not a commutative diagram). 5 marks on the page. NODES (5): (1) an
%   unlabelled endpoint at upper left, the far end of the leftmost ray; (2) $s_2$ at the top centre, the
%   far end of the vertical ray; (3) $s_1$ at the upper right, the far end of the right ray; (4) the
%   ellipsis "..." printed between the leftmost ray and the vertical ray, roughly at mid-height,
%   slightly left of the vertical; (5) $0$ at the bottom centre, the common origin of the three rays,
%   its label "0" printed to the right of the vertex. EDGES (3), all plain straight undirected line
%   segments (no arrowheads, no arrowtips), all emanating from the vertex $0$: (a) $0$
%   up-and-to-the-left to the unlabelled endpoint (1); no label; (b) $0$ straight up to $s_2$; labelled
%   $c_2$, the label printed to the RIGHT of the segment, about two-thirds of the way up; (c) $0$
%   up-and-to-the-right to $s_1$; labelled $c_1$, the label printed BELOW-RIGHT of the segment, about
%   halfway along. No arrows are hooked, dashed, mono, epi or iso; every stroke is a plain solid
%   segment.
Au voisinage de $s \in S$ et du point singulier $x_s$ de $H_s \subset X$ , on se trouve (dans des coordonnées locales convenables) dans la situation étudié en XIV 3.2. Si $\alpha$ est choisi holomorphe près de $s$ , on peut même trouver des coordonnées locales sur $X$ et $\ell$ telles que $f$ s'écrive $\Sigma z_i^2$ et que $\alpha^{-1}(c_s)$ soit " $z$ réel positif" . Au-dessus de l'intersection d'un voisinage de $s$ et de $\alpha^{-1}(c_s) - s$ , on peut alors considérer le fibré en sphères réelles $\Delta'_s \subset X$ , défini par " $z_i$ réel" . Pour $t \to s$ , la fibre $\Delta'_{st}$ de $\Delta'_s$ en $t$ tend vers $x_s$ . On notera $\Delta_s$ la réunion de $\{x_s\}$ et d'un fibré en sphères au-dessus de $\alpha^{-1}(c_s) - \{s\}$ , contenu dans $X - T(A)$ et qui prolonge $\Delta'_s$ . $\Delta_s$ est un disque plongé dans $X$ , de bord $H_o \cap \Delta_s$ une sphère qui représente le cycle évanescent $\delta_s$ en $s$ , transporté dans $H_o$ le long du chemin $\alpha^{-1}(c_s)$ .

Appliquons la théorie de Morse à la fonction $|\alpha f|^2$ , pour étudier comment change l'homotopie de

\begin{equation*}
\{x \in X| \ x \in T(A) \quad \mathrm{ou} \quad |\alpha f(x)|^2 \leq r \} \quad 
\end{equation*}

%%%%%%%%%%  scan idx 344  |  folio 337  %%%%%%%%%%

 lorsque $r$ croît. Soit

\begin{equation*}
X^- = \{x \in X \mid x \in T(A) \quad \mathrm{ou} \quad \alpha f(x) \neq \infty \} \ .
\end{equation*}
On trouve que l'inclusion

\begin{equation*}
H_o \cup \bigcup_{s \in S} \Delta_s \hookrightarrow T(A) \cup H_o \cup \bigcup_{s \in S} \Delta_s \hookrightarrow X^-
\end{equation*}
est une équivalence d'homotopie. On sait, encore par la théorie de Morse, que $X - X^-$ peut se rétracter par déformation sur un sous-complexe de dimension réelle $\leq n$ . Au total, on a donc le résultat suivant

\underline{Théorème} 4.2. \quad \underline{Il existe un sous-complexe} $\Sigma$ \underline{compact et de dimension (réelle)} $\leq n$ \underline{de} $H_\infty^{\mathrm{aff}} = H_\infty - A$ \underline{tel que l'inclusion}

\begin{equation*}
H_o \cup \bigcup_{s \in S} \Delta_s \longrightarrow X - \Sigma
\end{equation*}
\underline{soit une équivalence d'homotopie.}

4.3 \quad La soustraction de $\Sigma$ est sans effet sur les groupes de cohomologie, d'homologie ou d'homotopie d'indice $i \leq n$ de $X$ . Puisque $H_o \cup \bigcup_{s \in S} \Delta_s$ se déduit de $H_o$ en attachant des $n+1$-disque le long de sphères représentant les cycles évanescents, on trouve en homologie

\begin{equation*}
\mathbb{Z}^S \longrightarrow H_n(V_o) \longrightarrow H_n(X) \longrightarrow 0 \ .
\end{equation*}
La suite exacte analogue en cohomologie, et celle qui s'en déduit par dualité de Poincaré, forment un diagramme commutatif

\begin{equation}\tag{4.3.1}
\begin{tikzcd}[row sep=large, column sep=large]
{} & \mathbb{Z}^S & {} & {} \\
\mathbb{Z}^S \arrow[r, "{\Sigma n_s \delta_s}"] & H^n(V_o) \arrow[u, "{( \ , \delta_s)}"'] \arrow[r] & H^{n+2}(X) \arrow[r] & 0 \\
{} & H^n(X) \arrow[u] \arrow[ur, "{\eta\wedge}"'] & {} & {} \\
{} & 0 \arrow[u] & {} &
{} \end{tikzcd}
\end{equation}
% STRUCTURE: Tag "(4.3.1)" printed at the far left margin, on the level of the main horizontal row.
%   NODES (7): (1) $\mathbb{Z}^S$ at top, sitting directly above $H^n(V_o)$; (2) $\mathbb{Z}^S$ at the
%   left end of the main horizontal row; (3) $H^n(V_o)$, second entry of the main row; (4) $H^{n+2}(X)$,
%   third entry of the main row; (5) $0$, right end of the main row; (6) $H^n(X)$, one row below
%   $H^n(V_o)$ and directly beneath it; (7) $0$ at the bottom, directly beneath $H^n(X)$. ARROWS (7),
%   all solid single-shafted arrows with a plain ">" head; none are hooked, dashed, double-headed, mono,
%   epi or iso: (a) $\mathbb{Z}^S$ (2) $\longrightarrow$ $H^n(V_o)$ (3), horizontal, pointing RIGHT,
%   labelled $\Sigma n_s \delta_s$ ABOVE the shaft; (b) $H^n(V_o)$ (3) $\longrightarrow$ $H^{n+2}(X)$
%   (4), horizontal, pointing RIGHT, unlabelled; (c) $H^{n+2}(X)$ (4) $\longrightarrow$ $0$ (5),
%   horizontal, pointing RIGHT, unlabelled; (d) $H^n(V_o)$ (3) $\longrightarrow$ $\mathbb{Z}^S$ (1),
%   vertical, pointing UP, labelled $( \ , \delta_s)$ printed to the RIGHT of the shaft — the first slot
%   of the pairing is left BLANK on the page; (e) $H^n(X)$ (6) $\longrightarrow$ $H^n(V_o)$ (3),
%   vertical, pointing UP, unlabelled; (f) $0$ (7) $\longrightarrow$ $H^n(X)$ (6), vertical, pointing
%   UP, unlabelled; (g) $H^n(X)$ (6) $\longrightarrow$ $H^{n+2}(X)$ (4), diagonal, pointing
%   UP-AND-RIGHT, labelled $\eta\wedge$ printed BELOW-RIGHT of the shaft.

%%%%%%%%%%  scan idx 345  |  folio 338  %%%%%%%%%%


5. \underline{Les cas d'exception au théorème de Noether}

Dans ce numéro, nous persistons à travailler sur $\mathbb{C}$ , en cohomologie entière.

5.1 \quad Soient $2n$, $d$ et $\underline{a} = (a_1,\ldots,a_d)$ $(d \geq 1, a_i \geq 2)$ tels que la cohomologie de dimension $2n$ des intersections complètes lisses $V$ de dimension $2n$ et de multidegré $\underline{a}$ soit purement de type $(n, n)$ . En d'autres termes, $n = 0$ ou $(2n; \underline{a})$ est égal à $(2n; 2)$, $(2n; 2, 2)$ ou $(2; 3)$ (XI 2.9 ).

Supposons que $V$ se déduise d'une intersection complète $X \subset \mathbb{P}^{2n+d}$ (resp. $X \subset \mathbb{P}^{2n+d+1}$) de dimension $2n + 1$ et de multidegré $(a_1,\ldots,a_{d-1})$ (resp. $(a_1,\ldots,a_d)$ ) en prenant une section par une hypersurface de degré $a_d$ (resp. une section hyperplane) appartenant à un pinceau de Lefschetz $\ell$ .

La \underline{partie évanescente} $\mathrm{Ev}(V)$ de $H^{2n}(V)$ , égale à la \underline{partie primitive} de cette cohomologie est, d'après 4.3.1, engendrée \underline{sur} $\mathbb{Z}$ par les cycles évanescents. C'est l'orthogonal de la puissance $\eta^n$ de la classe de cohomologie $\eta$ d'une section hyperplane.

Puisque sur $H^{2n}(V)$ la forme d'intersection est unimodulaire, et que $\eta^{2n}$ n'est pas le multiple $by$ $(b > 1)$ d'une autre classe de cohomologie (XI 1.6 (iv)), le discriminant de la forme d'intersection sur $\mathrm{Ev}(V)$ est égal à

\begin{equation*}
(\eta^n, \eta^n) = \Pi\, a_i \quad .
\end{equation*}
Soit $R$ le système de racines dans $H^{2n}(V)$ formé des cycles évanescents et de leurs conjugués. On a sur $R$ les informations suivantes

(a) \quad $R$ est vide ou de type $A$, $D$ ou $E$ ;

(b) \quad le rang de $R$ est donné en XI 1.10 ;

(c) \quad le discriminant est $\Pi\, a_i$ . Dans les tables, ce discriminant est l'ordre du centre du groupe simple simplement connexe complexe correspondant. L'argument par lequel nous avons calculé le discriminant montre que ce centre, dual du quotient du réseau des poids par celui des racines, est cyclique.


%%%%%%%%%%  scan idx 346  |  folio 339  %%%%%%%%%%

Dans tous les cas de figure, ces informations déterminent $R$ complètement si $n \neq 0$ . On trouve:

\underline{Proposition} 5.2 \quad (i) \underline{Pour} $V$ \underline{uhe quadrique de dimension paire}, $R$ \underline{est de type} $A_1$

(ii) \underline{Pour} $V$ \underline{une surface cubique}, $R$ \underline{est de type} $E_6$ .

(iii) \underline{Pour} $V$ \underline{de dimension} $2n$ , \underline{et l'intersection de deux quadriques}, $R$ \underline{est de type} $D_{2n+3}$ .

5.3. \quad Supposons que $n \neq 0$ . L'ensemble $R^1$ des vecteurs de longueur $\pm 2$ dans $\mathrm{Ev}(V)$ est aussi un système de racine, ayant les mêmes caractéristiques (a) (b) (c) que $R$ . On a donc $R = R^1$ . Soit $\pi_1$ le groupe fondamental de l'espace des paramètres de toutes les intersections complètes $V \subset \mathbb{P}^{2n+d}$ du type considéré. L'image de $\pi_1$ dans $\mathrm{GL}(\mathrm{Ev}(V))$ est le groupe de Weil de $R$ . Dans tous les cas de figure, c'est le groupe d'automorphismes du triple

\begin{equation*}
(H^{2n}(V),\ \eta^n\ ,\ \mbox{forme d'intersection}).
\end{equation*}
5.4. \quad Avec $R$ remplacé par $R^1$ de 5.3, le résultat 4.2 (ii) est bien connu [1] et pour les petites valeurs de $n$ , 4.2 (iii) était connu, au moins par Manin ($n = 0$ est trivial; une surface intersection de deux quadriques est une surface de del Pezzo). Une étude détaillée de l'intersection de deux quadriques a été faite par M. Reid [4]. Il déduit 4.2 (iii) d'une étude de la géométrie des sous-espaces linéaires de dimension $n$ sur une intersection de deux quadriques $V \subset \mathbb{P}^{2n+2}$ .


\underline{Bibliographie}

[1] \quad E. CARTAN. Quelques remarques sur les 28 bitangentes d'une quartique plane et les 27 droites d'une surface cubique.\\ Bull. Sc. Math. \underline{70} 1946 p. 42-45, Oeuvres, partie I, tome 2 p. 1353.

[2] \quad P. DELIGNE. Théorie de Hodge II. Publ. Math. IHES \underline{40}, 1971 p. 5-58.


%%%%%%%%%%  scan idx 347  |  folio 340  %%%%%%%%%%

[3]\quad S. LEFSCHETZ. L'analysis situs et la géométrie algébrique. Gauthier-Villars 1924 (réédité dans Selected papers - Chelsea).

[4]\quad M. REID. Thèse, Cambridge 1972.

[5]\quad J.P. SERRE. Faisceaux algébriques cohérents. Ann. of Math. LXI (1955) p. 197-278 (cité FAC)
